% Vietnamese translation for OI Public Library
% Source-File: IOI中国国家候选队论文/2005/龙凡/龙凡.doc
\documentclass[11pt]{article}
\usepackage{oipl-vn}

\newcommand{\Oh}{\mathcal{O}}

\title{Ứng dụng của thứ tự}
\author{Long Fan\\Trường Trung học Yali, Trường Sa}
\date{2005}

\begin{document}
\maketitle

\origtitle{Ứng dụng của thứ tự}
\sourcepath{IOI China candidate team papers / 2005 / Long Fan.doc}

\renewcommand{\abstractname}{Tóm tắt}
\begin{abstract}
Bản chất của thi tin học là khai phá dữ liệu, còn ``thứ tự'' là một quan hệ
thường gặp nhưng khó phát hiện ẩn giữa dữ liệu. Nếu trong quá trình giải bài
ta tìm được quan hệ thứ tự ẩn trong đề, bài toán thường trở nên sáng rõ. Bài
viết tập trung thảo luận cách phát hiện và xây dựng thứ tự thích hợp trong các
quan hệ dữ liệu phức tạp, từ đó làm đơn giản và giải quyết bài toán.
\end{abstract}

\noindent\textbf{Từ khóa:} quan hệ dữ liệu; thứ tự; cây; dãy DFS; đồ thị; thứ
tự tô-pô.

\section{Dẫn nhập}

Thứ tự được dùng rất rộng trong tin học. Những ví dụ cơ bản nhất là tìm kiếm
nhị phân dựa trên thứ tự lớn nhỏ, và quy hoạch động trên đồ thị có hướng không
chu trình dựa trên thứ tự tô-pô. Để có những thứ tự ấy, ta có các thuật toán
tương ứng như sắp xếp nhanh và sắp xếp tô-pô.

Đồ thị, cây, dãy tuyến tính, tập hợp, v.v. là những quan hệ dữ liệu khác nhau,
và mỗi loại quan hệ có những thứ tự tương ứng. Ngay cả cùng một quan hệ dữ
liệu cũng có nhiều thứ tự khác nhau tùy theo ý nghĩa đang xét. Chẳng hạn, đồ
thị có hướng có thứ tự DFS theo nghĩa duyệt sâu, nhưng có thứ tự tô-pô theo
nghĩa phụ thuộc trước sau. Bản thân thứ tự không nhất thiết phải là tuyến
tính; thứ tự tô-pô cũng không phải một thứ tự tuyến tính chặt.

Trong dữ liệu rối rắm, tìm được thứ tự có giá trị và dùng nó hợp lý chính là
chủ đề của bài viết này.

\section{Dùng thứ tự để giải bài toán}

Nhiều bài tương tác sẽ trở nên dễ hơn nếu tìm được thứ tự thích hợp rồi khai
thác tính chất của nó. Xét ví dụ sau.

\subsection{Bài toán lưới}

\paragraph{Mô tả bài toán.}
Có một lưới điểm $N\times N$, với
\[
  1\le N\le 2003.
\]
Giữa hai điểm kề nhau có một cung có hướng mang trọng số nguyên $w$,
$1\le w\le 500000$. Đồng thời, từ điểm góc trên trái $(1,1)$ tới cùng một
điểm bất kỳ, mọi đường đi đều có cùng độ dài, tính bằng tổng trọng số các cung
đi qua.

Mỗi lần ta có thể hỏi hai giá trị liên quan tới một điểm $(x,y)$: khoảng cách
từ điểm đó tới điểm ngoài cùng bên phải cùng hàng, và khoảng cách từ điểm đó
tới điểm dưới cùng cùng cột. Cần tìm một điểm sao cho khoảng cách từ góc trên
trái tới điểm ấy bằng số nguyên cho trước $L$,
\[
  1\le L\le 2\cdot 10^9.
\]
Ban đầu chỉ biết $N$ và $L$, và được hỏi nhiều nhất $6667$ lần.

\paragraph{Phân tích.}
Trước hết có thể phủ định ý tưởng đơn giản nhất: hỏi để tính ra độ dài từng
cạnh rồi giải hệ phương trình. Tổng số cạnh là $2N(N-1)$. Từ điều kiện mọi
đường đi từ $(1,1)$ tới cùng một điểm có độ dài như nhau, ta chỉ có thể suy ra
khoảng $(N-1)^2$ phương trình độc lập, nên vẫn cần quá nhiều câu hỏi. Khi
$N=2003$, giới hạn $6667$ chỉ xấp xỉ $3N$, vì thế hướng này không thể dùng.

Hãy nhìn lại yêu cầu của đề: ta chỉ cần tìm một điểm có khoảng cách từ góc
trên trái bằng $L$. Ta không quan tâm trọng số từng cạnh, mà chỉ quan tâm
khoảng cách của từng điểm tới góc trên trái. Gọi giá trị đó là trọng số của
điểm.

Một sự thật quan trọng là: trọng số của một điểm luôn nhỏ hơn trọng số của mọi
điểm nằm phía dưới bên phải nó. Đây là thứ tự cơ bản nhất của bài toán. Nó
không dễ dùng trực tiếp, nhưng gợi ý rằng nếu trọng số tăng dần từ trái trên
sang phải dưới, thì trên các đường chéo từ trái dưới lên phải trên, giá trị
có thể được khai thác theo một cách khác.

Ý tưởng là tìm trước một điểm ở phía trái dưới có trọng số gần $L$, rồi đi dần
theo hướng phải trên. Cụ thể:
\begin{enumerate}
  \item Trên cột ngoài cùng bên trái, dùng tìm kiếm nhị phân để tìm điểm có
  trọng số vừa nhỏ hơn $L$, còn điểm ngay dưới nó đã lớn hơn $L$.
  \item Từ điểm đó, trên cùng hàng, dùng tìm kiếm nhị phân để tìm điểm gần
  $L$ nhất. Gọi điểm bắt đầu kiểm tra là $S$.
\end{enumerate}

Theo cách chọn $S$ và thứ tự cơ bản ở trên, các vùng bên trái, bên dưới, và
một phần bên phải dưới của $S$ đều không thể chứa nghiệm. Nghiệm nếu tồn tại
chỉ còn nằm trong vùng phía phải trên của điểm đang kiểm tra.

Khi xét điểm hiện tại:
\begin{itemize}
  \item Nếu trọng số bằng $L$, ta đã tìm được nghiệm.
  \item Nếu trọng số lớn hơn $L$, mọi điểm cùng hàng ở bên phải càng lớn hơn,
  nên không thể là nghiệm; ta dịch điểm kiểm tra lên trên.
  \item Nếu trọng số nhỏ hơn $L$, mọi điểm cùng cột ở phía trên không phù hợp
  theo vùng còn lại; ta dịch điểm kiểm tra sang phải.
\end{itemize}

Ta luôn giữ bất biến rằng nghiệm chỉ có thể nằm trong vùng phía phải trên của
điểm đang kiểm tra. Mỗi bước thu hẹp vùng ấy bằng cách đi lên hoặc sang phải.
Hai lần tìm kiếm nhị phân ban đầu cộng với quá trình đi đơn điệu này cần số
câu hỏi tuyến tính theo $N$, nhỏ hơn giới hạn khi $N$ đạt cực đại.

Từ một thứ tự rất đơn giản và cơ bản, cộng với việc thiết kế thuật toán phù
hợp với đặc điểm của nó, ta đã giải được bài toán. Điều này nhắc ta không nên
xem nhẹ những tính chất nội tại có vẻ đơn giản trong đề bài.

\section{Dùng thứ tự để đơn giản hóa bài toán}

\subsection{Bài toán dựng cây}

\paragraph{Mô tả bài toán.}
Cho một cây có $n$ đỉnh, đánh số $1,2,\ldots,n$. Mỗi đỉnh $v$ có một trọng số
$s(v)$, cũng là một hoán vị của $1,2,\ldots,n$. Với đỉnh $v$, định nghĩa
$t(v)$ là số hậu duệ của $v$ có trọng số nhỏ hơn $s(v)$. Biết cây và các giá
trị
\[
  t(1),t(2),\ldots,t(n),
\]
hãy tìm một cách gán trọng số cho các đỉnh. Nếu có nhiều nghiệm, in một nghiệm
bất kỳ.

\paragraph{Phân tích.}
Thoạt nhìn bài toán khó bắt đầu, nhưng chỉ cần nắm định nghĩa của $t(v)$ thì
cách giải trở nên nhẹ hơn. Định nghĩa ấy nói tới các hậu duệ có trọng số nhỏ
hơn $s(v)$, nên gợi ý dùng thứ tự duyệt trước DFS. Một tính chất quan trọng
của duyệt trước là: các hậu duệ của một đỉnh xuất hiện ngay sau đỉnh ấy trong
dãy duyệt, dưới dạng các đoạn con liên tiếp của các cây con.

Hãy tưởng tượng có $n$ ô trống từ trái sang phải, tương ứng với các trọng số
$1,2,\ldots,n$. Nếu ô thứ $j$ được điền đỉnh $i$, ta hiểu rằng $s(i)=j$.
Theo thứ tự duyệt trước DFS, lần lượt đặt từng đỉnh vào các ô. Khi xử lý đỉnh
$i$, ta đặt nó vào ô trống thứ $t(i)+1$ tính từ trái sang. Lý do là sau nó còn
có đúng $t(i)$ hậu duệ phải mang trọng số nhỏ hơn nó, tức phải được đặt vào
các ô trống bên trái nó. Nếu lần nào cũng chừa đúng $t(i)$ ô trống bên trái,
điều kiện của bài được thỏa mãn.

\paragraph{Chứng minh phác thảo.}
Trước hết chứng minh bằng quy nạp rằng với một cây con có $N$ đỉnh, thuật toán
sẽ đặt toàn bộ đỉnh của cây con vào $N$ ô trống đầu tiên đang dành cho cây con
đó. Với $N=1$ hiển nhiên đúng vì $t(v)=0$. Với $N=k$, gọi gốc là $r$ và các
cây con của nó lần lượt là $S_1,S_2,\ldots,S_m$. Dãy duyệt trước có dạng
\[
  r,\ S_1,\ S_2,\ \ldots,\ S_m.
\]
Sau khi đặt $r$ vào ô trống thứ $t(r)+1$, từng cây con có kích thước nhỏ hơn
$k$ được xử lý theo giả thiết quy nạp, nên toàn bộ cây con của $r$ cùng với
$r$ chiếm đúng các ô trống đầu tiên dành cho cây ấy.

Tiếp theo, với mỗi đỉnh $i$, các ô trống bên trái được chừa lại trước khi đặt
$i$ đều sẽ được lấp bởi các hậu duệ của chính nó. Điều này suy ra từ nhận xét
trên: đỉnh $i$ và các hậu duệ của nó chiếm một đoạn ô đầu tiên trong phạm vi
cây con của $i$.

Như vậy bài toán được đưa về dạng: có một dãy $N$ ô trống, mỗi lần nhận lệnh
``đặt số $A_i$ vào ô trống thứ $P_i$'', cần tìm trạng thái cuối cùng. Bài toán
này có thể giải bằng cây Fenwick hoặc cây đoạn: ghi nhận ô nào còn trống, và
mỗi lần tìm ô trống thứ $P_i$ trong $\Oh(\log N)$. Tổng độ phức tạp là
\[
  \Oh(N\log N),
\]
phần DFS chỉ mất $\Oh(N)$, còn bộ nhớ là $\Oh(N)$.

\subsection{Bài toán xếp hàng binh sĩ}

\paragraph{Mô tả bài toán.}
Có $N$ binh sĩ và một hàng vị trí đánh số từ trái sang phải. Mỗi lệnh của
tướng có dạng \texttt{Goto(L,S)}.
\begin{itemize}
  \item Nếu vị trí $L$ đang trống, binh sĩ $S$ đứng vào vị trí $L$.
  \item Nếu vị trí $L$ đang có binh sĩ $K$, thì $S$ đứng vào vị trí $L$, còn
  $K$ bị đẩy sang bằng cách thực hiện tiếp \texttt{Goto(L+1,K)}.
\end{itemize}
Tướng ban ra $N$ lệnh; mỗi binh sĩ được nhắc tới đúng một lần. Cần tìm trạng
thái hàng cuối cùng. Trong quá trình thực hiện, vị trí binh sĩ đứng có thể
vượt quá giới hạn ban đầu, nên trước hết cần xác định độ dài cuối cùng của
hàng; ô trống được in là $0$.

\paragraph{Phân tích ban đầu.}
Mô phỏng trực tiếp có thể mất $\Oh(N^2)$ vì một lần chèn gây ra dây chuyền
đẩy dài. Có thể tối ưu mã bằng cấu trúc hợp nhất tập hợp và thao tác dịch
khối, nhưng về bản chất độ phức tạp xấu nhất vẫn chưa thay đổi.

Một cách mô phỏng khác là xét cả khối liên tiếp. Một khối là một đoạn binh sĩ
liền nhau. Khối được xem như một chỉnh thể đối với bên ngoài, còn bên trong
lưu vị trí tương đối của các thành viên. Vì cần tìm kiếm và chèn phần tử trong
khối, ta dễ nghĩ tới cây cân bằng, trong đó khóa của mỗi phần tử là vị trí
tương đối bên trong khối. Tuy nhiên lại xuất hiện vấn đề gộp khối: nếu khoảng
cách giữa hai khối trở thành $0$, chúng phải hợp nhất. Cây cân bằng không dễ
hợp nhất; cách thô là tháo từng phần tử của khối nhỏ để chèn vào khối lớn.
Phân tích khấu hao cho thời gian xấu nhất khoảng $\Oh(N\log^2 N)$, và cài đặt
khá rườm rà. Một cách dùng cây cân bằng tốt hơn là xây dựng cây trên toàn bộ
các ô theo thứ tự, mỗi nút ghi số ô trống để định vị và xóa; cách này đạt
$\Oh(N\log N)$.

Bài toán còn có một cách đơn giản và hiệu quả hơn. Điểm then chốt là thao tác
chèn gây ra dây chuyền đẩy. Xét ví dụ có hai lệnh
\[
  \texttt{Goto(2,1)},\quad \texttt{Goto(2,2)}.
\]
Khi thực hiện lệnh thứ hai, binh sĩ $1$ bị đẩy. Nhưng nếu ta điền binh sĩ
$2$ trước, rồi điền binh sĩ $1$ vào ô trống thứ hai, hiệu quả cuối cùng là
tương đương. Điều này giống hệt bài toán dựng cây sau khi chuyển về dạng điền
số vào ô trống.

Định nghĩa lệnh mới \texttt{newgoto(L,S)}: đặt binh sĩ $S$ vào ô trống thứ
$L$, trong đó $L$ là thứ tự trong các ô trống chứ không phải chỉ số vị trí
tuyệt đối. Câu hỏi trở thành: với một dãy lệnh \texttt{Goto}, có tồn tại một
dãy lệnh \texttt{newgoto} tương đương chỉ bằng cách đổi thứ tự các cặp
$(L,S)$ hay không?

Câu trả lời là có. Thứ tự cần tìm phải thỏa hai điều kiện:
\begin{enumerate}
  \item Nếu hai khối không liên quan $A,B$ và $A$ nằm bên trái $B$, thì mọi
  phần tử của $B$ phải xuất hiện trước mọi phần tử của $A$ trong dãy
  \texttt{newgoto}. Nếu chèn $A$ trước, nó sẽ làm vị trí của $B$ dịch sang
  phải.
  \item Nếu việc chèn binh sĩ $S$ làm binh sĩ $K$ bị đẩy, thì $S$ phải được
  chèn trước $K$ trong dãy \texttt{newgoto}. Đây chính là nguyên lý tránh dây
  chuyền đẩy.
\end{enumerate}
Nếu xem mỗi binh sĩ là một đỉnh, và đặt cạnh có hướng $A\to B$ khi $A$ phải
được chèn trước $B$, thì thứ tự cần tìm là một thứ tự tô-pô của đồ thị này.

Ta không thể xây dựng đồ thị rồi sắp xếp tô-pô trực tiếp: việc tìm từng cạnh
gần như đã là mô phỏng lệnh \texttt{Goto}, và số cạnh có thể là bậc hai. Thay
vào đó, dùng mô phỏng không hoàn chỉnh: khi mô phỏng, không lưu quan hệ vị trí
bên trong khối, vì quan hệ ấy không ảnh hưởng tới thứ tự cần tìm.

Với mỗi khối, dùng cấu trúc hợp nhất tập hợp để lưu thông tin cơ bản của khối
và một dãy lệnh \texttt{newgoto} riêng của khối, ký hiệu $\operatorname{cmd}(A)$.
Ban đầu mọi vị trí đều trống. Khi xử lý lệnh \texttt{Goto(\(L_i\),i)}:

\paragraph{Chèn vào ô trống.}
Nếu $i$ được chèn vào một vị trí trống, ta tạo khối mới. Dãy lệnh của khối này
chỉ gồm chính cặp lệnh
\[
  \operatorname{cmd}(i)=(L_i,i).
\]

\paragraph{Chèn vào một khối.}
Nếu $i$ được chèn vào một khối $A$, ta đặt $i$ vào khối và cập nhật dãy lệnh
bằng cách đưa cặp $(L_i,i)$ lên đầu:
\[
  \operatorname{cmd}(A+i)=(L_i,i)+\operatorname{cmd}(A).
\]
Phần bị đẩy do $i$ chèn vào phải xuất hiện sau $i$; phần không bị đẩy cũng
phải xuất hiện sau $i$ để vị trí của $i$ không bị ảnh hưởng. Vì vậy đưa
$(L_i,i)$ lên đầu là hợp lý.

\paragraph{Hợp nhất khối.}
Một lần chèn có thể làm hai khối liền nhau và cần hợp nhất. Nếu khối $A$ nằm
bên trái khối $B$, thì mọi phần tử của $B$ phải xuất hiện trước mọi phần tử
của $A$ trong dãy \texttt{newgoto}. Do đó
\[
  \operatorname{cmd}(A+B)=\operatorname{cmd}(B)+\operatorname{cmd}(A).
\]
Các trường hợp hợp nhất thường gặp đều xử lý theo nguyên tắc này.

Sau khi mô phỏng hết, lấy các khối còn lại theo thứ tự từ phải sang trái và
nối các dãy $\operatorname{cmd}$ của chúng lại. Dãy thu được chính là thứ tự
tô-pô cần tìm, và dãy \texttt{newgoto} này tương đương với dãy \texttt{Goto}
ban đầu.

Quá trình biến đổi lệnh có thể làm gần tuyến tính bằng hợp nhất tập hợp và
danh sách nối. Sau đó mô phỏng dãy \texttt{newgoto} bằng cây Fenwick hoặc cây
đoạn, giống bài dựng cây, trong $\Oh(N\log N)$. Tổng độ phức tạp là
\[
  \Oh(N\log N).
\]

Qua hai bài toán trên, ta đã dùng những thứ tự khác nhau để đưa các bài tưởng
như không liên quan về cùng một dạng: điền số vào ô trống thứ $k$, giải bằng
cây Fenwick hoặc cây đoạn. Thứ tự thích hợp giúp ta thấy bản chất của bài
toán rõ hơn.

\section{Kết luận}

Những bài toán khác nhau ẩn chứa những thứ tự khác nhau để ta khai phá. Nếu
biết tận dụng, thứ tự sẽ như một đôi mắt giúp ta nhìn thấu bài toán hơn, và
nhìn rõ bản chất trong những quan hệ dữ liệu rối rắm.

Các ví dụ trên cũng cho thấy việc dùng thứ tự không hề đơn giản. Có những thứ
tự ẩn rất sâu, thậm chí phải tự xây dựng mới có. Điều này đòi hỏi tích lũy
thường xuyên và sự nhạy cảm với các quan hệ tinh tế giữa dữ liệu.

\section*{Lời cảm ơn}

Xin cảm ơn bạn Lei Tao vì bài gốc tự sáng tác và bạn Li Shi vì phần cải biên
đề bài. Xin cảm ơn Ren Kai và Zhou Yuan đã giúp đỡ trong bài toán xếp hàng
binh sĩ.

\section*{Tài liệu tham khảo}

\begin{itemize}
  \item Đề thi CEOI 2003.
  \item \textit{Introduction to Algorithms}.
  \item Đề gốc do Lei Tao, Trường Trung học Yali, Trường Sa, Hồ Nam, sáng tác.
  \item Luận văn đội tuyển tập huấn quốc gia năm 2002 của Li Rui.
\end{itemize}

\end{document}
